\subsection{Nikodym Metric and Setwise Measure Limits}

The Nikodym metric turns measurable sets modulo $\mu$-null symmetric difference into a complete metric space. This lets Baire category control setwise limits of absolutely continuous measures.

\begin{definition} \label{def:setwise-convergence-and-uniform-absolute-continuity} \hlt{Setwise Convergence and Uniform Absolute Continuity}\\
Let $(X,\mathcal{M})$ be a measurable space.
\begin{enumerate}[label=(\roman*)]
\setlength{\itemsep}{0pt}
\item For $A,B\in\mathcal{M}$, $A\triangle B\defeq(A\setminus B)\cup(B\setminus A)$.
\item Finite measures $\nu_n$ \hlt{converge setwise} to $\nu:\mathcal{M}\to[0,+\infty]$ iff $\nu_n(E)\to\nu(E)$ for every $E\in\mathcal{M}$.
\item For $\nu_n\ll\mu$ for all $n$, $(\nu_n)$ is \hlt{uniformly absolutely continuous} w.r.t. $\mu$ iff\\
$\forall\varepsilon>0\ \exists\delta>0$: $\mu(E)<\delta\Rightarrow\sup\limits_n\nu_n(E)<\varepsilon$ for all $E\in\mathcal{M}$.
\end{enumerate}
\end{definition}

\begin{lemma} \label{lem:nikodym-symmetric-difference-calculus} \hlt{Nikodym Symmetric-Difference Calculus}\\
For $A,B,C\in\mathcal{M}$ and a finite measure $\nu$:
\begin{enumerate}[label=(\roman*)]
\setlength{\itemsep}{0pt}
\item $A\triangle C\subseteq(A\triangle B)\cup(B\triangle C)$.
\item $\abs{\nu(A)-\nu(B)}\leq\nu(A\triangle B)$.
\end{enumerate}
\begin{proof}
\begin{enumerate}[label=(\roman*)]
\setlength{\itemsep}{0pt}
\item If $x\in A\triangle C$, then $\mathbf{1}_A(x)\neq\mathbf{1}_C(x)$. Hence $\mathbf{1}_A(x)\neq\mathbf{1}_B(x)$ or $\mathbf{1}_B(x)\neq\mathbf{1}_C(x)$, so $x\in(A\triangle B)\cup(B\triangle C)$.
\item Since $\nu(A)-\nu(B)=\nu(A\setminus B)-\nu(B\setminus A)$, $\abs{\nu(A)-\nu(B)}\leq\nu(A\setminus B)+\nu(B\setminus A)=\nu(A\triangle B)$.
\end{enumerate}
\hlt{Structural Commentary}: Set inclusion gives the triangle inequality; disjoint differences give the mass bound.\\
\hlt{Key Insight}: Nikodym distance measures exactly the part where two measurable sets disagree.
\end{proof}
\end{lemma}

\begin{definition} \label{def:nikodym-metric-space} \hlt{Nikodym Metric Space}\\
Let $(X,\mathcal{M},\mu)$ be a finite measure space.
\begin{enumerate}[label=(\roman*)]
\setlength{\itemsep}{0pt}
\item \hlt{Null equivalence}: $A\sim_\mu B\Leftrightarrow\mu(A\triangle B)=0$.
\item \hlt{Quotient space}: $\mathcal{M}_\mu\defeq\mathcal{M}/{\sim_\mu}$, with class notation $[A]\in\mathcal{M}_\mu$.
\item \hlt{Metric}: $\rho_\mu([A],[B])\defeq\mu(A\triangle B)$.
\end{enumerate}
By Lemma \ref{lem:nikodym-symmetric-difference-calculus} and Proposition \ref{prop:pseudometric-metric-quotient}, $\rho_\mu$ is a well-defined metric on $\mathcal{M}_\mu$.
\end{definition}

\begin{proposition} \label{prop:nikodym-pseudometric} \hlt{Nikodym Pseudometric}\\
Let $E\subseteq\R$ be Lebesgue measurable with $m(E)<\infty$, let $X$ be the collection of Lebesgue measurable subsets of $E$, and for $A,B\in X$ set $A\Delta B\defeq(A\setminus B)\cup(B\setminus A)$ and $\rho(A,B)\defeq m(A\Delta B)$. Define $A\sim_\rho B$ iff $\rho(A,B)=0$.
\begin{enumerate}[label=(\roman*)]
\setlength{\itemsep}{0pt}
\item $\rho$ is a pseudometric on $X$.
\item $\widetilde{\rho}([A],[B])\defeq\rho(A,B)$ is a metric on $X/{\sim_\rho}$.
\item For all $A,B\in X$, $\rho(A,B)=\int_E\abs{\mathbf{1}_A-\mathbf{1}_B}\,dm$.
\end{enumerate}
\begin{proof}
\begin{enumerate}[label=(\roman*)]
\setlength{\itemsep}{0pt}
\item For $A,B,C\in X$, $A\Delta B\subseteq E$, $A\Delta A=\emptyset$, $A\Delta B=B\Delta A$, and $A\Delta C\subseteq(A\Delta B)\cup(B\Delta C)$.\\
Hence $0\leq\rho(A,B)\leq m(E)<\infty$, $\rho(A,A)=0$, $\rho(A,B)=\rho(B,A)$, and $\rho(A,C)\leq\rho(A,B)+\rho(B,C)$.
\item By (i), $\rho$ is a pseudometric and $A\sim_\rho B\iff\rho(A,B)=0$.\\
Proposition \ref{prop:pseudometric-metric-quotient} gives that $\widetilde{\rho}([A],[B])\defeq\rho(A,B)$ is a metric on $X/{\sim_\rho}$.
\item For $t\in E$, $\abs{\mathbf{1}_A(t)-\mathbf{1}_B(t)}=\mathbf{1}_{A\Delta B}(t)$; integrating gives $\int_E\abs{\mathbf{1}_A-\mathbf{1}_B}\,dm=m(A\Delta B)=\rho(A,B)$.
\end{enumerate}
\hlt{Structural Commentary}: Symmetric difference measures set disagreement.\\
\hlt{Key Insight}: Quotienting null disagreement gives definiteness.
\end{proof}
\end{proposition}

\begin{theorem} \label{thm:nikodym-metric-complete} \hlt{Completeness of the Nikodym Metric}\\
If $(X,\mathcal{M},\mu)$ is a finite measure space, then $(\mathcal{M}_\mu,\rho_\mu)$ is complete.
\begin{proof} Let $(A_n)$ be $\rho_\mu$-Cauchy. Choose $n_k\uparrow\infty$ with $\rho_\mu(A_{n_{k+1}},A_{n_k})<2^{-k}$, and set $D_k\defeq A_{n_{k+1}}\triangle A_{n_k}$.\\
Since $\mu(\bigcup\limits_{k\geq N}D_k)\leq\sum\limits_{k\geq N}2^{-k}\to0$, $N_0\defeq\bigcap\limits_N\bigcup\limits_{k\geq N}D_k$ is null. Put $A\defeq\bigcup\limits_N\bigcap\limits_{k\geq N}A_{n_k}$.\\
Then $A_{n_k}\triangle A\subseteq N_0\cup\bigcup\limits_{j\geq k}D_j$, so $\rho_\mu(A_{n_k},A)\leq\sum\limits_{j\geq k}2^{-j}\to0$. Since $(A_n)$ is Cauchy, $\rho_\mu(A_n,A)\to0$.\\
\hlt{Structural Commentary}: Summable jumps give an a.e. eventual set limit.\\
\hlt{Key Insight}: Summable symmetric differences confine infinitely changing membership to a null set.
\end{proof}
\end{theorem}

\begin{proposition} \label{prop:absolute-continuity-nikodym-uniform-continuity} \hlt{Absolutely Continuous Measures Are Nikodym-Continuous}\\
Let $(X,\mathcal{M},\mu)$ be finite, let $\nu$ be a finite measure with $\nu\ll\mu$, and define $\Phi_\nu([E])\defeq\nu(E)$.
\begin{enumerate}[label=(\roman*)]
\setlength{\itemsep}{0pt}
\item $\Phi_\nu:\mathcal{M}_\mu\to\R$ is well-defined.
\item $\Phi_\nu$ is uniformly continuous.
\end{enumerate}
\begin{proof}
\begin{enumerate}[label=(\roman*)]
\setlength{\itemsep}{0pt}
\item If $A\sim_\mu B$, then $\mu(A\triangle B)=0$, hence $\nu(A\triangle B)=0$ and $\nu(A)=\nu(B)$.
\item By Lemma \ref{lem:nikodym-symmetric-difference-calculus}, $\abs{\nu(A)-\nu(B)}\leq\nu(A\triangle B)$. Given $\varepsilon>0$, choose $\delta>0$ by Lemma \ref{lem:absolute-continuity-epsilon-delta-finite-source} such that $\mu(E)<\delta$ implies $\nu(E)<\varepsilon$. Thus $\rho_\mu(A,B)<\delta$ implies $\abs{\Phi_\nu(A)-\Phi_\nu(B)}<\varepsilon$.
\end{enumerate}
\hlt{Structural Commentary}: Null-set identification defines the map; absolute continuity controls its modulus.\\
\hlt{Key Insight}: Absolute continuity converts small Nikodym distance into small change of mass.
\end{proof}
\end{proposition}

\begin{proposition} \label{prop:uniform-absolute-continuity-rn-derivatives} \hlt{Uniform Absolute Continuity and Radon--Nikodym Derivatives}\\
Let $(X,\mathcal{M},\mu)$ be finite, let $(\nu_n)$ be finite measures with $\nu_n\ll\mu$, and set $f_n=\frac{d\nu_n}{d\mu}$. The following are equivalent.
\begin{enumerate}[label=(\roman*)]
\setlength{\itemsep}{0pt}
\item $(\nu_n)$ is uniformly absolutely continuous with respect to $\mu$.
\item The functions $\Phi_{\nu_n}:\mathcal{M}_\mu\to\R$ are uniformly equicontinuous.
\item The functions $\Phi_{\nu_n}$ are equicontinuous at $\emptyset$.
\item $\{f_n:n\geq1\}$ is uniformly integrable: $\forall\varepsilon>0\ \exists\delta>0\ \forall E\in\mathcal{M}$, $\mu(E)<\delta\Rightarrow\sup\limits_n\int_Ef_n\,d\mu<\varepsilon$.
\end{enumerate}
\begin{proof}
\begin{enumerate}[label=(\roman*)]
\setlength{\itemsep}{0pt}
\item[$(i)$]$\Rightarrow(ii)$ Given $\varepsilon>0$, choose $\delta>0$ with $\mu(E)<\delta\Rightarrow\sup\limits_n\nu_n(E)<\varepsilon$.\\
If $\rho_\mu(A,B)=\mu(A\triangle B)<\delta$, then $\abs{\Phi_{\nu_n}(A)-\Phi_{\nu_n}(B)}=\abs{\nu_n(A)-\nu_n(B)}\leq\nu_n(A\triangle B)<\varepsilon$ uniformly in $n$.
\item[$(ii)$]$\Rightarrow(iii)$ Given $\varepsilon>0$, uniform equicontinuity gives $\delta>0$ such that\\
$\rho_\mu(A,B)<\delta\Rightarrow\sup\limits_n\abs{\Phi_{\nu_n}(A)-\Phi_{\nu_n}(B)}<\varepsilon$. With $B=\emptyset$, this is equicontinuity at $\emptyset$.
\item[$(iii)$]$\Rightarrow(i)$ Given $\varepsilon>0$, choose $\delta>0$ such that $\rho_\mu(E,\emptyset)<\delta\Rightarrow\sup\limits_n\abs{\Phi_{\nu_n}(E)-\Phi_{\nu_n}(\emptyset)}<\varepsilon$.\\
Since $\rho_\mu(E,\emptyset)=\mu(E)$ and $\Phi_{\nu_n}(\emptyset)=0$, $\mu(E)<\delta\Rightarrow\sup\limits_n\nu_n(E)<\varepsilon$.
\item[$(i)$]$\Leftrightarrow(iv)$ For $E\in\mathcal{M}$, $\nu_n(E)=\int_Ef_n\,d\mu$. Thus $\sup\limits_n\nu_n(E)<\varepsilon\Leftrightarrow\sup\limits_n\int_Ef_n\,d\mu<\varepsilon$ for every small-$\mu$ set $E$.
\end{enumerate}
\hlt{Structural Commentary}: Uniform AC of $\nu_n$ equals uniform integrability of $f_n$.\\
\hlt{Key Insight}: Radon--Nikodym turns set bounds into integral bounds.
\end{proof}
\end{proposition}

\begin{theorem} \label{thm:setwise-limit-under-uniform-absolute-continuity} \hlt{Setwise Limits under Uniform Absolute Continuity}\\
Let $(X,\mathcal{M},\mu)$ be finite, let $\nu_n$ be finite measures with $\nu_n\ll_{\mathrm{u}}\mu$, $\sup\limits_n\nu_n(X)<+\infty$, and $\nu_n(E)\to\nu(E)$ $\forall E\in\mathcal{M}$.
\begin{enumerate}[label=(\roman*)]
\setlength{\itemsep}{0pt}
\item $\nu$ is a finite measure.
\item $\nu\ll\mu$.
\end{enumerate}
\begin{proof}
\begin{enumerate}[label=(\roman*)]
\setlength{\itemsep}{0pt}
\item Since $\nu(X)=\lim\limits_n\nu_n(X)\leq\sup\limits_n\nu_n(X)$, $\nu$ is finite-valued. Finite additivity passes to the pointwise limit.\\
For pairwise disjoint $(E_k)$, set $E\defeq\bigcup\limits_kE_k$ and $R_N\defeq\bigcup\limits_{k>N}E_k$.\\
Given $\varepsilon>0$, uniform absolute continuity and $\mu(R_N)\downarrow0$ give $\nu_n(R_N)<\varepsilon$ $\forall n$ and all sufficiently large $N$, hence $\nu(R_N)\leq\varepsilon$. Therefore $\nu(E)=\sum\limits_{k=1}^N\nu(E_k)+\nu(R_N)\to\sum\limits_k\nu(E_k)$.
\item If $\mu(A)=0$, then for every $\varepsilon>0$, uniform absolute continuity gives $\sup\limits_n\nu_n(A)<\varepsilon$.\\
Passing to the setwise limit, $\nu(A)\leq\varepsilon$ $\forall\varepsilon>0$, hence $\nu(A)=0$.
\end{enumerate}
\hlt{Structural Commentary}: Uniform tail control yields countable additivity and preserves the null ideal.\\
\hlt{Key Insight}: Uniform absolute continuity makes the tail of any disjoint union uniformly small.
\end{proof}
\end{theorem}

\begin{lemma} \label{lem:nikodym-one-point-equicontinuity-spreads} \hlt{One-Point Equicontinuity Spreads}\\
Let $(X,\mathcal{M},\mu)$ be finite, and let $(\nu_n)$ be finite measures with $\nu_n\ll\mu$.\\
If the maps $\Phi_{\nu_n}$ are equicontinuous at one point $E_0\in\mathcal{M}_\mu$, then they are uniformly equicontinuous on $\mathcal{M}_\mu$.
\begin{proof}
Fix $\varepsilon>0$ and choose $\delta>0$ so that $\rho_\mu(E,E_0)<\delta$ implies $\sup\limits_n\abs{\nu_n(E)-\nu_n(E_0)}<\frac{\varepsilon}{2}$.\\
If $\mu(C)<\delta$, then both $E_0\setminus C$ and $E_0\cup(C\setminus E_0)$ are within $\delta$ of $E_0$, so $\nu_n(C\cap E_0)<\frac{\varepsilon}{2},\; \nu_n(C\setminus E_0)<\frac{\varepsilon}{2}$ $\forall n$.\\
Hence $\sup\limits_n\nu_n(C)<\varepsilon$. Taking $C=A\triangle B$ gives $\sup\limits_n\abs{\nu_n(A)-\nu_n(B)}<\varepsilon$ whenever $\rho_\mu(A,B)<\delta$.\\
\hlt{Structural Commentary}: Split a small set into its portions inside and outside the equicontinuity base point.\\
\hlt{Key Insight}: Smallness near one set controls both the part inside it and the part outside it.
\end{proof}
\end{lemma}

\begin{figure}[H]
  \centering
  \begin{tikzpicture}[
  >=Stealth,
  every node/.style={font=\scriptsize},
  setbox/.style={draw=blue!60!black, fill=blue!5, rounded corners=1.5pt,
    minimum width=2.35cm, minimum height=0.66cm, align=center,
    inner xsep=5pt, inner ysep=4pt},
  measbox/.style={draw=green!45!black, fill=green!8, rounded corners=1.5pt,
    minimum width=2.35cm, minimum height=0.66cm, align=center,
    inner xsep=5pt, inner ysep=4pt},
  catbox/.style={draw=red!55!black, fill=red!5, rounded corners=1.5pt,
    minimum width=2.35cm, minimum height=0.66cm, align=center,
    inner xsep=5pt, inner ysep=4pt},
  arrow/.style={->, line width=0.55pt, black!65},
  dasharrow/.style={->, dashed, line width=0.5pt, black!55},
  edgeword/.style={font=\tiny, fill=white, inner sep=1.2pt, align=center}
]
  \node[setbox, minimum width=2.70cm] (space) at (-2.55,1.35)
    {$\mathcal{M}_\mu$ complete\\Nikodym space};

  \node[measbox] (measures) at (-6.00,0)
    {$\nu_n\ll\mu$\\$\nu_n\to\nu$ setwise};
  \node[measbox, minimum width=2.70cm] (functions) at (-2.45,0)
    {$\Phi_n([E])=\nu_n(E)$\\[-1pt] continuous\\[-1pt] pointwise limit};
  \node[catbox] (baire) at (1.10,0)
    {Baire point\\of equicontinuity};
  \node[measbox] (uniform) at (4.50,0)
    {uniform absolute\\continuity};
  \node[measbox] (limit) at (7.85,0)
    {$\nu$ is a measure\\and $\nu\ll\mu$};

  \draw[arrow] (measures) -- (functions);
  \draw[arrow] (functions) -- (baire);
  \draw[arrow] (baire) -- (uniform);
  \draw[arrow] (uniform) -- (limit);
  \draw[dasharrow] (space.south east) -- node[right, edgeword] {Baire} (baire.north);
\end{tikzpicture}

  \caption{Vitali--Hahn--Saks mechanism: setwise convergence becomes pointwise convergence on $\mathcal{M}_\mu$; Baire gives uniform absolute continuity.}
  \label{fig:nikodym-vitali-hahn-saks-flow}
\end{figure}

\begin{method} \label{meth:vitali-hahn-saks-baire-mechanism} \hlt{Baire Mechanism for Vitali--Hahn--Saks}\\
Regard each absolutely continuous measure $\nu_n$ as $\Phi_{\nu_n}:\mathcal{M}_\mu\to\R$, $\Phi_{\nu_n}(E)\defeq\nu_n(E)$.
\begin{enumerate}[label=(\roman*)]
\setlength{\itemsep}{0pt}
\item \hlt{Translate}: Setwise convergence becomes pointwise convergence of $(\Phi_{\nu_n})$.
\item \hlt{Apply Baire}: Completeness of $\mathcal{M}_\mu$ supplies a point of equicontinuity.
\item \hlt{Spread}: Lemma \ref{lem:nikodym-one-point-equicontinuity-spreads} upgrades one-point equicontinuity to uniform equicontinuity.
\item \hlt{Return to measures}: Proposition \ref{prop:uniform-absolute-continuity-rn-derivatives} converts uniform equicontinuity into uniform absolute continuity.
\end{enumerate}
\end{method}

\begin{theorem} \label{thm:vitali-hahn-saks-setwise-convergence} \hlt{Vitali--Hahn--Saks Theorem}\\
Let $(X,\mathcal{M},\mu)$ be finite, and let $(\nu_n)$ be finite measures with $\nu_n\ll\mu$, $\sup\limits_n\nu_n(X)<+\infty$, and $\nu_n\to\nu$ setwise.
\begin{enumerate}[label=(\roman*)]
\setlength{\itemsep}{0pt}
\item $(\nu_n)$ is uniformly absolutely continuous with respect to $\mu$.
\item $\nu$ is a finite measure.
\item $\nu\ll\mu$.
\end{enumerate}
\begin{proof}
\begin{enumerate}[label=(\roman*)]
\setlength{\itemsep}{0pt}
\item Set $\Phi_n([E])\defeq\nu_n(E)$. By Proposition \ref{prop:absolute-continuity-nikodym-uniform-continuity}, each $\Phi_n$ is continuous on complete $\mathcal{M}_\mu$.\\
If $A\sim_\mu B$, then $\nu_n(A)=\nu_n(B)$ $\forall n$, hence $\nu(A)=\nu(B)$.\\
Thus $\Phi([E])\defeq\nu(E)$ is well-defined and $\Phi_n\to\Phi$ pointwise. Since $\abs{\Phi_n([E])}\leq\sup\limits_n\nu_n(X)$ and $\nu_n(E)\to\nu(E)$, $\Phi$ is real-valued.\\
Theorem \ref{thm:pointwise-limit-continuous-on-dense-set} gives an equicontinuity point $E_0\in\mathcal{M}_\mu$ for $(\Phi_n)$.\\
Lemma \ref{lem:nikodym-one-point-equicontinuity-spreads} gives uniform equicontinuity on $\mathcal{M}_\mu$, and Proposition \ref{prop:uniform-absolute-continuity-rn-derivatives} gives $\nu_n\ll_{\mathrm{u}}\mu$.
\item By (i), $\nu_n\ll_{\mathrm{u}}\mu$; Theorem \ref{thm:setwise-limit-under-uniform-absolute-continuity}(i) gives that $\nu$ is a finite measure.
\item By (i), $\nu_n\ll_{\mathrm{u}}\mu$; Theorem \ref{thm:setwise-limit-under-uniform-absolute-continuity}(ii) gives $\nu\ll\mu$.
\end{enumerate}
\hlt{Structural Commentary}: Baire gives one equicontinuity point; Nikodym calculus spreads it.\\
\hlt{Key Insight}: Setwise convergence plus bounded total mass forces uniform null-set control.
\end{proof}
\end{theorem}

\begin{corollary} \label{cor:nikodym-bounded-setwise-limits} \hlt{Nikodym Bounded Setwise Limit Theorem}\\
If finite measures $(\nu_n)$ on $(X,\mathcal{M})$ converge setwise to $\nu$ and $\sup\limits_n\nu_n(X)<+\infty$, then $\nu$ is finite.
\begin{proof}
Define
$\mu(E)\defeq\sum\limits_{n=1}^{\infty}2^{-n}\nu_n(E).$
Then $\mu$ is a finite measure by countable additivity of nonnegative series, and $\nu_n\ll\mu$ $\forall n$ because $2^{-n}\nu_n\leq\mu$. Thus $\nu$ is a finite measure by Theorem \ref{thm:vitali-hahn-saks-setwise-convergence}(ii).\\
\hlt{Structural Commentary}: A summable mixture places the sequence under one finite dominating measure.\\
\hlt{Key Insight}: A summable mixture builds one finite measure that dominates the whole sequence.
\end{proof}
\end{corollary}
